
\subsection*{Model covariates}

Variables based on land-use \cite{hudson2017}, population density \cite{gpw2017}, and road network density \cite{groads2013} data, formed the core of all models in the study\cite{depalma2021}. For land-use, the land-use class and its intensity was combined into a new set of categorical variables; records where this information was unknown were filtered out. Minimally used plantation forest was grouped with lightly used secondary forest, and other plantation forest (light and intense use) was grouped with intensely used secondary forest \cite{depalma2021}. The categorical land-use variables were one-hot encoded, with the model intercepts representing minimally used primary vegetation reference sites. Human population density and road network density at 1 and 50~km$^2$ scales, log transformed to reduce skewness, were also part of all models. The BTMs additionally included the following bioclimatic and topographic variables: annual mean precipitation and temperature (1~km$^2$), temperature and precipitation seasonality (1~km$^2$), elevation (1 and 10~km$^2$), and terrain roughness (1 and 10~km$^2$). A complete list of model covariates can be found in \autoref{tab:variables}. We did not include any interaction effects due to the small number of observations in many studies (and to some extent, hierarchical groups), implying that the models were already highly parameterized, and even overparameterized to some extent, given the variables above (see \autoref{fig:sites_per_study} for data coverage and below for a discussion on implications of this).

To derive the continuous variables, the coordinates of each sampling site were first projected from global EPSG:4326 to local UTM format, before generating circular polygons of different spatial extents. The equal-area projections ensured that the calculated values were comparable across all locations, regardless of distance to the equator. For each raster dataset (human population density, bioclimatic, topographic) the mean value of each polygon was calculated (including partially covered pixels), after which the polygons were reprojected to the global format. For the road network data, a similar approach was used to derive road density as the combined length of all roads within each polygon. Population density data were interpolated between the available years in the GPW dataset, and back to the earliest PREDICTS data (1984), assuming an exponential growth rate \cite{goldewijk2005}. The population densities were matched to the sampling year of each site; the other covariates were static layers.

For the beta diversity (BC) models, we used the alpha diversity model covariates as a starting point. Here, the land-use variables described the conditions at the non-reference site in each pair, since the reference sites constitute the model intercepts (minimally used primary vegetation). Differences in the human population and road density variables were calculated for each site pair. Additionally, the spatial (Haversine) distance, and multivariate environmental (Gower) distance \cite{gower1971}, were included to control for natural compositional turnover \cite{depalma2021}. The Haversine distance was normalized by the median sampling extent among all studies. The Gower distance was calculated using the following variables, all at the 1~km$^2$ scale: maximum and minimum temperature of the warmest and coldest months, respectively, precipitation of the wettest and driest months, and elevation. All continuous variables were standardized prior to model training, and evaluated for collinearity within their respective category.
\vspace{\baselineskip}
